\documentclass[11pt]{amsart}
\usepackage[margin=1.1in]{geometry}
\usepackage{amsmath,amssymb,amsthm,mathtools,enumitem}
\usepackage[colorlinks=true,linkcolor=blue!50!black,citecolor=blue!50!black,urlcolor=blue!50!black]{hyperref}

\newtheorem{theorem}{Theorem}
\newtheorem{proposition}[theorem]{Proposition}
\newtheorem{lemma}[theorem]{Lemma}
\newtheorem{corollary}[theorem]{Corollary}
\theoremstyle{definition}
\newtheorem{definition}[theorem]{Definition}
\theoremstyle{remark}
\newtheorem{remark}[theorem]{Remark}
\newtheorem{question}[theorem]{Question}

\newcommand{\R}{\mathbb{R}}
\newcommand{\Z}{\mathbb{Z}}
\newcommand{\dH}{d_{\mathrm H}}

\title{Thin thickenings of the plane need seven colors}
\author{Ansh Mishra}
\author{Aryan Senthilkumar}
\date{\today}

\begin{document}

\begin{abstract}
For a set $X\subset\R^d$ let $\chi(X)$ be the least number of colors in a coloring of $X$
with no two points at Euclidean distance exactly one having the same color. We show that for
every $\delta>0$ there is a finite set $K$ of diameter less than $\delta$ in some Euclidean
space such that $\chi(\R^2\times K)=7$. Consequently $\chi(\R^2\times[0,\varepsilon]^k)=7$
whenever $\varepsilon$ is small and $k\ge C/\varepsilon$ for an absolute constant $C$. This
answers, with $k$ depending on $\varepsilon$, a question of Voronov, who asked whether
$\chi(\R^2\times[0,\varepsilon]^k)=7$ for small $\varepsilon$ and some $k$ independent of
$\varepsilon$. The colorings are arbitrary: no measurability or regularity is assumed. The
proof combines Voronov's theorem on colorings of the plane with a forbidden interval of
distances and the Frankl--R\"odl theorem on forbidden distances in codes, in the cross form
of Keevash and Long.
\end{abstract}

\maketitle

\section{Introduction}

The Hadwiger--Nelson problem asks for the chromatic number $\chi(\R^2)$ of the graph on the
plane in which two points are adjacent when their distance is exactly one. The hexagonal
coloring gives $\chi(\R^2)\le 7$. De~Grey~\cite{deGrey} proved $\chi(\R^2)\ge5$, and the
value is still open.

Several relaxations are better understood. For $0<\varepsilon<1$ write
$\chi_{[1-\varepsilon,1+\varepsilon]}(\R^2)$ for the least number of colors in a coloring of
the plane in which no two points at distance in $[1-\varepsilon,1+\varepsilon]$ share a color.
Voronov~\cite{Voronov} proved:

\begin{theorem}[Voronov~{\cite[Theorem~1.1]{Voronov}}]\label{thm:voronov}
For every $\varepsilon>0$, $\chi_{[1-\varepsilon,1+\varepsilon]}(\R^2)\ge7$.
\end{theorem}

A different relaxation thickens the plane in extra dimensions while keeping the single
forbidden distance one. It is known that $\chi(\R^2\times[0,\varepsilon]^2)\ge 6$ and
$\chi(\R^3\times[0,\varepsilon]^6)\ge10$ for small $\varepsilon$
(see~\cite[Section~7]{Voronov} and the references there), and Voronov asked:

\begin{question}[{\cite[Question~7.2]{Voronov}}]\label{q:voronov}
Is it true that $\chi(\R^2\times[0,\varepsilon]^k)=7$ for sufficiently small $\varepsilon>0$
and some $k$ independent of $\varepsilon$?
\end{question}

Our main result says that seven colors are needed, and suffice, once the number of extra
dimensions is allowed to grow like $1/\varepsilon$.

\begin{theorem}\label{thm:main}
For every $\delta>0$ there exist $n\in\mathbb N$ and a finite set $K\subset\R^{2n}$ with
$\operatorname{diam}K<\delta$ such that
\[
  \chi(\R^2\times K)=7 .
\]
\end{theorem}

\begin{corollary}\label{cor:slab}
There are absolute constants $C,\varepsilon_0>0$ such that
$\chi(\R^2\times[0,\varepsilon]^k)=7$ whenever $0<\varepsilon<\varepsilon_0$ and
$C/\varepsilon\le k\le \varepsilon^{-2}/10$.
\end{corollary}

Here $\R^2\times K$ is a union of finitely many parallel translates of a plane in
$\R^{2+2n}$, all within distance $\delta$ of one another. The upper bound
$\chi\le7$ is elementary (Proposition~\ref{prop:upper}); the content is the lower bound,
which holds for arbitrary colorings.

The lower bound is a special case of a general transference principle.

\begin{theorem}\label{thm:transfer}
Let $m\ge1$ and $c\ge1$ be integers, and suppose that for every $\varepsilon>0$ every
$c$-coloring of $\R^m$ has two points of the same color at distance in
$[1-\varepsilon,1+\varepsilon]$. Then for every $\delta>0$ there is a finite set $K$ of
diameter less than $\delta$ in some Euclidean space such that $\chi(\R^m\times K)\ge c+1$.
\end{theorem}

Theorem~\ref{thm:main} follows from Theorem~\ref{thm:transfer} with $m=2$, $c=6$ and
Theorem~\ref{thm:voronov}, together with Proposition~\ref{prop:upper}.

\subsection*{Idea of the proof}
A coloring of $\R^m\times K$ assigns to each point $p$ of a fine lattice in $\R^m$ a coloring of
the fiber $\{p\}\times K$. We take $K$ to be a scaled image of the ternary Hamming cube
$\{0,1,2\}^n$, embedded so that squared Euclidean distances are proportional to Hamming
distances. Let $m(p)$ be the most frequent color on the fiber over $p$. If two lattice points
$p,p'$ whose distance is slightly less than one receive the same majority color, the two
majority classes are dense subsets of $\{0,1,2\}^n$; by the Frankl--R\"odl theorem for codes,
in the cross form of Keevash and Long, they contain a pair at any prescribed Hamming distance in
the bulk range. Choosing that distance so that the lift of the pair has Euclidean length exactly
one gives a monochromatic unit pair. Hence the majority coloring of the lattice avoids a whole
interval of distances below one, and after extending it to the plane this contradicts the
forbidden-interval hypothesis.

\section{Tools}

For $s,t\in\{0,1,2\}^n$ let $\dH(s,t)$ be the number of coordinates in which they differ.

\begin{theorem}[Keevash--Long~{\cite[Corollary~9]{KeevashLong}}; Frankl--R\"odl~\cite{FranklRodl}]
\label{thm:KL}
For every $\eta\in(0,1/2)$ and $q\ge3$ there is $\gamma=\gamma(\eta,q)>0$ such that the
following holds for every $n$. If $C,D\subset[q]^n$ satisfy $|C||D|>(q-\gamma)^{2n}$, then for
every integer $h\in(\eta n,(1-\eta)n)$ there are $x\in C$ and $y\in D$ with $\dH(x,y)=h$.
\end{theorem}

For $q=3$, $\eta=1/4$, $\gamma=\gamma(1/4,3)$ and an integer $c\ge1$, let $n_0(c)$ be the
least integer with $\bigl(3/(3-\gamma)\bigr)^{2n_0(c)}>c^2$. Then for $n\ge n_0(c)$ any two sets
$C,D\subset\{0,1,2\}^n$ with $|C|,|D|\ge 3^n/c$ satisfy $|C||D|\ge 3^{2n}/c^2>(3-\gamma)^{2n}$,
the hypothesis of Theorem~\ref{thm:KL}.

\medskip
\noindent\emph{The embedding.} Let $\tau(0)=(0,0)$, $\tau(1)=(1,0)$,
$\tau(2)=(1/2,\sqrt3/2)$, the vertices of a unit equilateral triangle. For $\lambda>0$ define
\[
  \iota_\lambda:\{0,1,2\}^n\to\R^{2n},\qquad
  \iota_\lambda(s)=\lambda\bigl(\tau(s_1),\dots,\tau(s_n)\bigr).
\]
Then
\begin{equation}\label{eq:iso}
  |\iota_\lambda(s)-\iota_\lambda(t)|^2=\lambda^2\,\dH(s,t)\qquad(s,t\in\{0,1,2\}^n),
\end{equation}
the image lies in $[0,\lambda]^{2n}$, and its diameter is $\lambda\sqrt n$.

\section{The transference theorem}

\begin{proof}[Proof of Theorem~\ref{thm:transfer}]
Fix $\delta>0$. Let $Q$ be a large positive integer, put $\lambda=1/Q$, and let
\[
  n=\max\Bigl(n_0(c),\ \Bigl\lceil \tfrac{12\sqrt m}{\lambda}\Bigr\rceil+1\Bigr),
  \qquad K=\iota_\lambda(\{0,1,2\}^n)\subset\R^{2n}.
\]
For $Q$ large we have $n\le 12\sqrt m\,Q+2$, so $\lambda^2n\le 12\sqrt m\,\lambda+2\lambda^2$
tends to zero. Choose $Q$ so large that $\operatorname{diam}K=\lambda\sqrt n<\delta$ and
$\tfrac34\lambda^2n<1$.

Suppose, for a contradiction, that $\phi:\R^m\times K\to\{1,\dots,c\}$ has no monochromatic
pair at distance one. Let $\Lambda=\lambda\Z^m$. For $p\in\Lambda$ define
$f_p(s)=\phi(p,\iota_\lambda(s))$ on $\{0,1,2\}^n$, let $\mu(p)$ be a most frequent value of
$f_p$, and put $A_p=f_p^{-1}(\mu(p))$, so that $|A_p|\ge3^n/c$.

\emph{Step 1: $\mu$ avoids an interval of lattice distances.} Let
\[
  I=\Bigl(\sqrt{1-\tfrac34\lambda^2n},\ \sqrt{1-\tfrac14\lambda^2n}\Bigr).
\]
Let $p,p'\in\Lambda$ with $|p-p'|\in I$. Since $|p-p'|^2=\lambda^2N$ with $N\in\Z$ and
$1=\lambda^2Q^2$, the number $h=Q^2-N$ is an integer with
$|p-p'|^2+\lambda^2h=1$, and $|p-p'|\in I$ means exactly $h\in(n/4,3n/4)$.
If $\mu(p)=\mu(p')=j$, then $|A_p|,|A_{p'}|\ge 3^n/c$ with $n\ge n_0(c)$, so
Theorem~\ref{thm:KL} gives $s\in A_p$ and $t\in A_{p'}$ with $\dH(s,t)=h$. By~\eqref{eq:iso},
\[
  \bigl|(p,\iota_\lambda(s))-(p',\iota_\lambda(t))\bigr|^2
  =|p-p'|^2+\lambda^2h=1,
\]
and both points have color $j$, a contradiction. Hence $\mu(p)\ne\mu(p')$ whenever
$|p-p'|\in I$.

\emph{Step 2: extension to $\R^m$.} For $x\in\R^m$ let $\pi(x)\in\Lambda$ be a nearest lattice
point (any fixed choice), so $|x-\pi(x)|\le\rho:=\lambda\sqrt m/2$, and set
$g(x)=\mu(\pi(x))$. Write $I=(a,b)$ and let $J=[a+3\rho,\,b-3\rho]$. Using
$\sqrt u-\sqrt v=(u-v)/(\sqrt u+\sqrt v)\ge(u-v)/2$ for $0\le v<u\le1$,
\[
  b-a\ \ge\ \tfrac14\lambda^2 n\ >\ \tfrac14\lambda^2\cdot\tfrac{12\sqrt m}{\lambda}
  \ =\ 3\sqrt m\,\lambda\ =\ 6\rho ,
\]
so $J$ is a nondegenerate interval. If $|x-y|\in J$ then
$|\pi(x)-\pi(y)|\in[a+\rho,\,b-\rho]\subset I$, hence $g(x)\ne g(y)$. Thus $g$ is a $c$-coloring of $\R^m$ with no monochromatic pair at distance
in $J=[\alpha,\beta]$, $0<\alpha<\beta$.

\emph{Step 3: rescaling.} The map $x\mapsto \tfrac{\alpha+\beta}{2}x$ transforms $g$ into a
$c$-coloring of $\R^m$ with no monochromatic pair at distance in
$[1-\varepsilon',1+\varepsilon']$, where $\varepsilon'=(\beta-\alpha)/(\beta+\alpha)>0$.
This contradicts the hypothesis. Therefore $\chi(\R^m\times K)\ge c+1$.
\end{proof}

\begin{remark}
The coloring $\phi$ in the proof is arbitrary; no measurability is used. Only finitely many
fibers matter for each unit pair, so by the de Bruijn--Erd\H{o}s theorem the conclusion is
witnessed by a finite unit-distance graph in $\R^m\times K$, although our argument does not
exhibit it.
\end{remark}

\section{The upper bound and the proofs of the main results}

\begin{proposition}\label{prop:upper}
If $K\subset\R^k$ has $\operatorname{diam}K\le 0.4$, then $\chi(\R^2\times K)\le7$.
\end{proposition}

\begin{proof}
Tile the plane by regular hexagons of circumradius $r=0.45$ and color them with seven colors
by the cosets of the index-seven sublattice of the hexagonal lattice, as in the classical
construction; boundary points receive the color of any incident hexagon. Two points in the
closure of one hexagon are at distance at most $2r=0.9$, and two points in the closures of
distinct hexagons of the same color are at distance at least $(\sqrt{21}-2)r>1.16$. Color
$(x,s)\in\R^2\times K$ by the color of $x$. If $|(x,s)-(y,t)|=1$, then
$|x-y|^2=1-|s-t|^2\ge 1-0.16$, so $0.91<|x-y|\le1$, and $x,y$ lie in hexagons of different
colors.
\end{proof}

\begin{proof}[Proof of Theorem~\ref{thm:main}]
Given $\delta>0$, apply Theorem~\ref{thm:transfer} with $m=2$, $c=6$ and $\min(\delta,0.4)$
in place of $\delta$; its hypothesis is Theorem~\ref{thm:voronov}. This gives $K$ with
$\operatorname{diam}K<\min(\delta,0.4)$ and $\chi(\R^2\times K)\ge7$, and
Proposition~\ref{prop:upper} gives equality.
\end{proof}

\begin{proof}[Proof of Corollary~\ref{cor:slab}]
In the proof of Theorem~\ref{thm:transfer} with $m=2$ the set $K=\iota_\lambda(\{0,1,2\}^n)$
lies in $[0,\lambda]^{2n}$, with $n=\max(n_0(6),\lceil 12\sqrt2/\lambda\rceil+1)$. Given small
$\varepsilon$, choose $\lambda=1/Q\le\varepsilon$ with $\lambda\ge\varepsilon/2$. Then
$\R^2\times K\subset\R^2\times[0,\varepsilon]^{2n}\subset\R^2\times[0,\varepsilon]^k$ for
every $k\ge2n$ (pad with zero coordinates), and $2n\le C/\varepsilon$ for a suitable absolute
$C$ once $\varepsilon$ is small, so $\chi(\R^2\times[0,\varepsilon]^k)\ge7$. If moreover
$k\le\varepsilon^{-2}/10$, the slab $[0,\varepsilon]^k$ has diameter
$\varepsilon\sqrt k<0.4$ and Proposition~\ref{prop:upper} gives $\chi\le7$.
\end{proof}

\section{Remarks}

\begin{remark}[Why $k$ grows]
In the argument the lattice spacing $\lambda$ is also the fiber scale, because lifted squared
distances $|p-p'|^2+\lambda^2h$ must hit the value one exactly. The interval avoided in
Step~1 has length of order $\lambda^2n$, and it must exceed the lattice spacing $\lambda$ for
Step~2. This forces $n\gtrsim1/\lambda$. To answer Question~\ref{q:voronov} as stated one
would need fibers of bounded dimension in which dense subsets realize every squared distance in
a window of width much larger than the lattice spacing, which requires number-theoretic input
(for instance distance theorems for dense subsets of $\{0,\dots,L\}^k$, $k\ge5$) together with
control of local obstructions.
\end{remark}

\begin{remark}[Other dimensions]
Theorem~\ref{thm:transfer} applies to every $m$. For $m=3$ it yields
$\chi(\R^3\times K)\ge10$ from the known bound $\chi_{[1-\varepsilon,1+\varepsilon]}(\R^3)\ge10$;
this is weaker than the known $\chi(\R^3\times[0,\varepsilon]^6)\ge10$, whose number of extra
dimensions does not depend on $\varepsilon$.
\end{remark}

\begin{remark}[Contrast with the plane]
For the plane itself the published bounds are $5\le\chi(\R^2)\le7$. By
Theorem~\ref{thm:main}, adding finitely many parallel copies of the plane at arbitrarily small
mutual distances already forces the value seven.
\end{remark}

\begin{thebibliography}{9}

\bibitem{deGrey}
A.~D.~N.~J. de~Grey, The chromatic number of the plane is at least 5,
\emph{Geombinatorics} \textbf{28} (2018), 18--31.

\bibitem{FranklRodl}
P.~Frankl and V.~R\"odl, Forbidden intersections,
\emph{Trans. Amer. Math. Soc.} \textbf{300} (1987), 259--286.

\bibitem{KeevashLong}
P.~Keevash and E.~Long, Frankl--R\"odl type theorems for codes and permutations,
\emph{Trans. Amer. Math. Soc.} \textbf{369} (2017), 1147--1162; arXiv:1402.6294.

\bibitem{Voronov}
V.~Voronov, The chromatic number of the plane with an interval of forbidden distances is at
least 7, arXiv:2304.10163 (version 3, 2025).

\end{thebibliography}

\end{document}
